\begin{table*}[t]
\centering
\small
\caption{\textbf{Instructions for human validation.}
Question and rating levels shown to annotators for each criterion, with our English translation.
The first level of each criterion is the passing level.
H5 is worded differently for A2 and A3 targets.}
\label{tab:human_rubric}
\begin{tabular}{@{}p{0.05\textwidth}p{0.44\textwidth}p{0.44\textwidth}@{}}
\toprule
 & \textbf{Korean} & \textbf{English translation} \\
\midrule
H1
& 한국어 가독성: 문장 자체가 자연스럽고 의미가 명확한가?
\newline \textsc{natural}: 한 번에 자연스럽게 읽힌다.
\newline \textsc{minor\_issue}: 조금 어색하지만 의미는 분명하다.
\newline \textsc{unnatural}: 어색함이나 모호함이 이해를 방해한다.
& Korean readability: Are the sentences themselves natural and is their meaning clear?
\newline \textsc{natural}: Reads naturally at once.
\newline \textsc{minor\_issue}: Slightly awkward, but the meaning is clear.
\newline \textsc{unnatural}: Awkwardness or ambiguity hinders understanding. \\
\midrule
H2
& 출제 인위성: 실제 상황처럼 자연스럽게 읽히는가, 아니면 정답과 선택지를 만들기 위해 조건을 인위적으로 나열한 글처럼 읽히는가?
\newline \textsc{natural}: 현실에서 있을 법한 자연스러운 상황으로 읽힌다.
\newline \textsc{noticeable}: 인위적인 구성 흔적은 있지만 하나의 상황으로 읽힌다.
\newline \textsc{contrived}: 실제 상황보다 문제 풀이를 위한 조건 나열처럼 읽힌다.
& Construction artifice: Does the text read naturally like a real situation, or like a list of conditions arranged to produce the answer and the options?
\newline \textsc{natural}: Reads as a natural situation that could occur in reality.
\newline \textsc{noticeable}: Shows traces of construction but still reads as a single situation.
\newline \textsc{contrived}: Reads as a list of conditions for solving the question rather than a real situation. \\
\midrule
H3
& 표면적 정답 누출: 전문 내용을 몰라도 길이·문법·형식만으로 정답을 찾을 수 있는가?
\newline \textsc{sound}: 표면 단서로 선택지를 줄이기 어렵다.
\newline \textsc{narrowing}: 일부는 지워지지만 둘 이상 남는다.
\newline \textsc{isolating}: 표면 단서만으로 정답 하나가 남는다.
& Surface answer leakage: Can the answer be found from length, grammar, or format alone, without knowing the subject matter?
\newline \textsc{sound}: The options are hard to narrow down with surface cues.
\newline \textsc{narrowing}: Some options are eliminated, but two or more remain.
\newline \textsc{isolating}: Surface cues alone leave only the gold answer. \\
\midrule
H4
& 정답 근거 추적: 정답을 anchor와 사례만으로 뒷받침할 수 있는가?
\newline \textsc{traceable}: 필요한 핵심 근거가 anchor에 있다.
\newline \textsc{partial}: 관련 근거는 있지만 anchor 밖 지식이 더 필요하다.
\newline \textsc{out\_of\_scope}: 핵심 근거를 anchor에서 찾을 수 없다.
& Traceability of the answer: Can the gold answer be supported by the anchor and the case alone?
\newline \textsc{traceable}: The key grounds needed are in the anchor.
\newline \textsc{partial}: Relevant grounds exist, but knowledge outside the anchor is also needed.
\newline \textsc{out\_of\_scope}: The key grounds cannot be found in the anchor. \\
\midrule
H5 (A2)
& 목표 구조: 정확히 하나의 mapping이 적용되고, 조건 충족을 사례에서 추론해야 하는가?
\newline \textsc{clear}: 하나의 mapping이 적용되고 조건 충족을 사례에서 분명히 추론해야 한다.
\newline \textsc{borderline}: 추론은 있으나 조건이 거의 드러나 있거나 mapping 수가 다소 불분명하다.
\newline \textsc{not\_met}: 직접 적용·조회이거나, 여러 mapping이 필요하거나, 사례로 결정할 수 없다.
& Target structure: Does exactly one mapping apply, and must the satisfaction of its condition be inferred from the case?
\newline \textsc{clear}: One mapping applies, and whether its condition is met must clearly be inferred from the case.
\newline \textsc{borderline}: Inference is involved, but the condition is nearly stated or the number of mappings is somewhat unclear.
\newline \textsc{not\_met}: The question is direct application or retrieval, requires several mappings, or cannot be decided from the case. \\
\midrule
H5 (A3)
& 목표 구조: 둘 이상의 mapping이 실제로 경쟁하고, 선택에 따라 정답이 달라지는가?
\newline \textsc{clear}: 둘 이상의 mapping이 사례에 근거하며, 지배 mapping에 따라 정답이 달라진다.
\newline \textsc{borderline}: 두 번째 mapping은 있으나 근거가 약하거나 정답이 달라지는지 불분명하다.
\newline \textsc{not\_met}: 실질적 경쟁이 없거나, mapping이 달라도 정답이 같거나, 결론을 정할 수 없다.
& Target structure: Do two or more mappings genuinely compete, and does the answer change with the choice between them?
\newline \textsc{clear}: Two or more mappings are grounded in the case, and the answer depends on which mapping governs.
\newline \textsc{borderline}: A second mapping exists, but its grounds are weak or it is unclear whether the answer would change.
\newline \textsc{not\_met}: There is no real competition, the answer is the same under either mapping, or the conclusion cannot be determined. \\
\bottomrule
\end{tabular}
\end{table*}
